\begin{theorem}\label{thm:equivalencias_conservativos}
Sea $\mathbf{F}: A \subseteq \mathbb{R}^n \to \mathbb{R}^n$ un campo vectorial continuo, con $A$ un conjunto abierto y arcoconexo. Las siguientes afirmaciones son equivalentes:
\begin{enumerate}
    \item $\mathbf{F}$ es un campo vectorial conservativo en $A$.
    \item La integral de línea de $\mathbf{F}$ es independiente del camino en $A$.
    \item $\oint_C \mathbf{F} \cdot d\mathbf{s} = 0$ para toda curva simple, cerrada y orientada $C$ contenida en $A$.
\end{enumerate}
\end{theorem}

    \textit{Demostración.} 
    \begin{itemize}
\item[] ($1 \Rightarrow 2.)$  Ver \ref{tmh:c2}
\item[] ($2 \Rightarrow 3.)$ Sea $C^{+}$ una curva simple, cerrada y orientada en $A$. Pensemos que a dicha curva la podemos ver como la uni\'on de  dos curvas   simples orientadas en $A$,  $C_1$ , $C_2$,  cuyas orientaciones son heredadas de $C$,   con los mismos puntos extremos,  tal como muestra la figura. 

\begin{figure}[htbp]
    \centering
    % Llamamos al archivo con el código TikZ extraído
    \begin{figure}[H]
        \begin{center}
            \begin{tikzpicture}[scale=2, >=stealth]

                % Puntos A y B
                \coordinate (A) at (0,0);
                \coordinate (B) at (3,0);

                % Primer camino: curva hacia arriba (gamma_1)
                \draw[thick, blue, ->, decoration={markings, mark=at position 0.5 with {\arrow{>}}},
                    postaction={decorate}] 
                    (A) .. controls (1,1.5) and (2,1.5) .. (B) node[midway, above] {$C_1$};

                % Segundo camino: curva hacia abajo (gamma_2)
                \draw[thick, red, ->, decoration={markings, mark=at position 0.5 with {\arrow{>}}},
                    postaction={decorate}] 
                    (B) .. controls (1,-1.5) and (2,-1.5) .. (A) node[midway, below] {$C_2$};

                % Nombres de los puntos
                \filldraw (A) circle(0.03) node[below left] {$A$};
                \filldraw (B) circle(0.03) node[below right] {$B$};

            \end{tikzpicture}
            \caption{Curvas $C_1$ y $C_2$.}
        \end{center}
    \end{figure} 
    \label{fig:campo conservativo}
\end{figure}

   Es decir,  $C^{+}=C_1^{+} \cup C_2^{-}$.  Así:
    \begin{align*}
        \oint_{C^{+}} \mathbf{F}\cdot d\mathbf{s}&=\int_{C_1^{+}}\mathbf{F}\cdot d\mathbf{s}+\int_{C_2^{-}}\mathbf{F}\cdot d\mathbf{s}\\
        &=\int_{C_1^{+}}\mathbf{F}\cdot d\mathbf{s}-\int_{C_2^{+}}\mathbf{F}\cdot d\mathbf{s} =0 \ 
          \end{align*}
   \item[] ($3 \Rightarrow 1$.) Queda a cargo del lector.   
\end{itemize}


\begin{tarea}  Probar que el campo $\mathbf{F}$  tiene rotor nulo en todo $\mathbb{R}^{2}$ y, sin embargo, no es un campo conservativo. ( No vale la vuelta de de la Propiedad (\ref{conserv})).   $$F(x,y) = (\dfrac{y}{x^{2}+y^{2}}, \dfrac{-x}{x^{2}+y^{2}}) $$
\end{tarea}



